% !TEX root = ../main.tex
\chapter{RESEARCH AND FINDINGS}
\label{chap:research}
% Does this chapter reuse your own published, presented, accepted, submitted,
% or preprint article? Add a chapter attribution here, directly below \label.
% It prints single spaced at the bottom of this chapter's first page, with no
% footnote number. Use the wording for the article's status (OGE manual 2026,
% section 2.8, Table 2.1), then the full citation in your citation style:
%   Published, presented, or a preprint:
%     Portions of this chapter have previously appeared as:
%     This chapter has previously appeared as:
%   Accepted, not yet published:
%     Portions of this chapter are to appear as:   (This chapter is to appear as:)
%   Submitted or under review (name the journal or conference):
%     Portions of this chapter are under review as:  (This chapter is ...)
%     Portions of this chapter have been submitted as:  (This chapter has ...)
% \fullcite{key} prints a full citation from references.bib. Delete the % signs:
% \chapterattribution{Portions of this chapter have previously appeared as:
%   \fullcite{your-article-key}.}
% Co-authored work also goes on the Declaration of Co-Contribution. If you are
% the sole author, keep the attribution; the declaration is not needed (OGE
% manual 2026, section 2.8.1).
Replace this chapter with your analysis, methods, findings, or argument.
The examples below are teaching examples, not research results.
Delete any examples you do not need.

\section{A table}
Table~\ref{tab:example} shows invented data to demonstrate a simple table.
Put the caption above a table. Use a full caption: the same words appear in the
list of tables automatically.

\begin{table}[htbp]
  \centering
  \caption{Illustrative values for three example groups.}
  \label{tab:example}
  % One letter per column (l, c, r: alignment); & next column, \\ ends a row.
  % Table too wide for the page? See "Figures and tables" in docs/writing.md.
  \begin{tabular}{lr}
    \toprule
    Group & Example value \\
    \midrule
    A & 12 \\
    B & 18 \\
    C & 24 \\
    \bottomrule
  \end{tabular}
\end{table}

\section{A figure}
Figure~\ref{fig:example} shows a sample image that comes with LaTeX. Upload
your own image to the figures folder (see figures/README.md).
Put figure captions below the image.

\begin{figure}[htbp]
  \centering
  % For your own image: \includegraphics[width=\linewidth]{figures/my-plot}
  % (the .png, .jpg, or .pdf extension is optional).
  \includegraphics[width=0.4\linewidth]{example-image}
  \caption{A sample image demonstrating the position of a figure caption.}
  \label{fig:example}
\end{figure}

\section{A reproduced figure}
Figure~\ref{fig:reproduced} shows where a permission statement goes when you
reuse someone else's figure or table. The bracketed text is a placeholder,
not a real source.

% A reproduced figure or table needs, in its caption: your description, the
% FULL citation of the original, and a permission or license statement.
% With permission (add any notice your permission document requires):
%   Reproduced from [full citation], with permission from [publisher].
% Creative Commons or open access (state the exact license):
%   Reproduced from [full citation] under the [CC BY 4.0] license.
% \fullcite{key} inside \caption{...} prints the full citation for you.
% A long caption is fine. Never use \caption[short]{long}: the list of
% figures must match the caption exactly.
% List your permission files in an appendix; see appendices/.
\begin{figure}[htbp]
  \centering
  \fbox{\parbox[c][1in][c]{0.8\linewidth}{\centering Reproduced figure goes here.}}
  \caption{An example caption for a reproduced figure. Reproduced from [full
    citation of the original source], with permission from [publisher].}
  \label{fig:reproduced}
\end{figure}

\section{An equation}
Delete this section if you do not need mathematics. Equation~\ref{eq:example}
demonstrates an equation with an automatically generated number:
\begin{equation}
  y = mx + b.
  \label{eq:example}
\end{equation}
Give every displayed equation a number by using the \verb|equation|
environment. Only displays within a lemma, algorithm, theorem, or proof may
be unnumbered. Equation~\ref{eq:multiline} spans two lines and has one centered number:
\begin{equation}
  \begin{split}
    (a + b)^2 &= (a + b)(a + b) \\
              &= a^2 + 2ab + b^2.
  \end{split}
  \label{eq:multiline}
\end{equation}
% If a number drops BELOW its equation, the line is too long. Break it with
% split as above (\\ ends a line, & aligns the lines) so the number stays on
% the same line, centered vertically. Use align for several numbered lines.
Use dollar signs for short mathematical expressions inside a sentence, such as
$x = 2$. Reference labels instead of typing page or equation numbers by hand.
